\subsection{Run 82: the entries obey Sato--Tate, and a fourth-moment
estimate}\label{sub:r82c}

Theorem \ref{thm:jacobi} turns every question about the family into a question
about the twisted Jacobi sums $J(\chi,\psi;q)$.  Two facts about them are
established here: their distribution, and a quantitative fourth-moment bound.

\emph{The distribution.}  Collapsing both character sums in
$\sum_{\chi,\psi\in S}|J|^{2k}$ over the index coset --- whose annihilator in
$\FF_N^*$ is $G_0=\langle a,-1\rangle$ of order $2K$ for odd $K$ --- imposes the
two conditions
\[
  y_1\cdots y_k=g\,x_1\cdots x_k,\qquad
  (x_1-1)\cdots(x_k-1)=h\,(y_1-1)\cdots(y_k-1),
  \qquad g,h\in G_0 .
\]
The diagonal solutions are the $k!$ matchings of $\{y\}$ with $\{x\}$ at
$g=h=1$, on which the phase $\sum x_i-\sum y_i$ vanishes identically, so the
naive prediction is $\langle|J|^{2k}\rangle=k!\,N^k$: complex-Gaussian moments.
Measured over five level sets at each of three designs:
\begin{center}\small
\begin{tabular}{rllll}
\toprule
$N$ & $\langle|J|^2\rangle/N$ & $\langle|J|^4\rangle/N^2$ & $\langle|J|^6\rangle/N^3$ & $\langle|J|^8\rangle/N^4$\\
\midrule
$199$ & $1.060623$ & $2.0143$ & $5.1119$ & $14.5205$\\
$3571$ & $1.007939$ & $2.0022$ & $5.0045$ & $14.0046$\\
$19489$ & $0.996038$ & $2.0062$ & $5.0361$ & $14.1649$\\
\midrule
Catalan $C_k$ & $1$ & $2$ & $5$ & $14$\\
Gaussian $k!$ & $1$ & $2$ & $6$ & $24$\\
\bottomrule
\end{tabular}
\end{center}
The moments are the Catalan numbers, with deviations decaying like $N^{-1/2}$.
Catalan numbers are the moments of the semicircle law on $[-2,2]$, so
\[
  J(\chi,\psi;q)/\sqrt N \;=\; \operatorname{tr}\theta,
  \qquad \theta\in SU(2)\ \text{equidistributed},
\]
the vertical Sato--Tate law of a \emph{rank-two} Kloosterman/hypergeometric
sheaf.  Equivalently: the naive diagonal count is correct for $k\le2$, where
$k!=C_k$, and from $k=3$ the \emph{crossing} matchings cancel --- the deficits
are $k!-C_k=0,0,1,10$ --- which is precisely the statement that the monodromy is
$SU(2)$.  This confirms, one level below the matrix elements, the
$SU(2)$-trace shape that Kurlberg and Rudnick conjecture for the matrix elements
of the desymmetrised cat map \cite{KR2003}.  It is also consistent with the GOE
statistics of \S\ref{sub:r82b}: universality for Wigner-type matrices requires
moment bounds, not Gaussianity.

\emph{A quantitative estimate.}  For $k=2$ the collapse is effective.  The two
constraints determine the elementary symmetric functions of $(y_1,y_2)$,
\[
  P:=y_1y_2=g\,x_1x_2,\quad
  Q:=(y_1-1)(y_2-1)=\frac{(x_1-1)(x_2-1)}{h},\quad
  S:=y_1+y_2=P+1-Q,
\]
so the fibre over $(x_1,x_2)$ has $1+\left(\tfrac{S^2-4P}{N}\right)$ points, and
$y_i\neq0,1$ automatically.  Hence
\[
  T_{g,h}(q)=\sum_{x_1,x_2\neq0,1}
    \Bigl(1+\Bigl(\tfrac{S^2-4P}{N}\Bigr)\Bigr)\,\ep\bigl(q(x_1+x_2-S)\bigr),
\]
a \emph{two}-variable complete sum, hence $O(N)$ by Weil--Deligne, where a naive
bound over the two-dimensional surface would give only $O(N^{3/2})$.  Since the
phase vanishes identically only on $g=h=1$ with $\{y\}=\{x\}$:

\begin{proposition}\label{prop:fourth}
$\displaystyle \frac1{m^2}\sum_{\chi,\psi\in S}|J(\chi,\psi;q)|^4
 = 2N^2\bigl(1+O(K^2/N)\bigr)$,
uniformly in $q\neq0$, the constant being that of the Weil--Deligne bound for
$T_{g,h}$.
\end{proposition}

The identity behind it is certified to $1.1\cdot10^{-14}$ over all $484$ tuples
at $N=199$; the diagonal term is exactly $q$-independent, with
$T_{1,1}/2N^2=0.9775$ at $N=199$ and $0.99874$ at $N=3571$; and
$\max_{(g,h)\neq(1,1)}|T_{g,h}(q)|/N$ equals $4.6$--$5.6$ at both designs, i.e.
bounded in $N$.  The measured relative error of the fourth moment is
$0.115/0.017/0.003$ at $N=199/3571/19489$, one factor of $K$ better than
Proposition \ref{prop:fourth}, reflecting square-root cancellation among the
$(2K)^2$ tuples.

Together with $\langle|J|^2\rangle=N+O(K^2)$ from the same collapse, this
supplies the first two entrywise moment estimates a Wigner-type local law
requires.  What it does not supply is the joint decorrelation of \emph{distinct}
entries with rates; that, as noted in \S\ref{sub:r82b}, is the single remaining
input, and it is a Katz-type equidistribution statement for the sheaf above.
